\documentclass[a4paper, 12pt]{article} %or book

\usepackage[english]{babel} %in English
\usepackage[utf8]{inputenc} %unicode input
\usepackage[T1]{fontenc}    %unicode output
\usepackage{listings}
\usepackage{graphicx}
\usepackage{float}

\usepackage{hyperref}       %hyperlinks
\hypersetup{colorlinks=true, %colour the links
	pdftitle="hyperlink example",linkcolor=blue,urlcolor=blue} %title of my PDF


\title{\textbf{Scaling Machine Learning with Spark}}	
\author{\textbf{Riya Pokhrel \& Lucas Torri}\\
	Université Bretagne Sud \& Paris Lodron Universität Salzburg}
\date{November 15, 2026}

\begin{document}
	\maketitle 
	
	\tableofcontents
	
	
	
	
	
	\section{Introduction}
	
	Airborne LiDAR is producing high resolution 3D point clouds of the territory. In France, the LiDAR HD programme of IGN (Institut National de l’information Geographique et forestière) aims to cover the whole territory by 2026. This means that this huge amount of point clouds have to be classified into land-cover classes. This volume cannot be processed into a single machine.
	
	The FRACTAL dataset (Gaydon et al., 2024) is a benchmark built from LiDAR HD data: 100,000 patches
	of $50 \times 50$~m labelled with seven classes and split into training, validation and test sets.
	Using deep neural network trained on GPUs, its authors reach 77.5\% mIoU.
	In this project, we use Apache Spark and MLlib on an AWS EMR cluster instead. Our objectives are to:
	\begin{itemize}
		\item build a PySpark pipeline that samples a fraction of the patches, computes features and
		trains a random forest;
		\item evaluate it with the per-class IoU, the mIoU and the overall accuracy;
		\item time the preprocessing, training and evaluation stages separately and draw speedup curves
		for several dataset fractions and numbers of executors;
		\item keep the experiments reproducible on larger fractions and clusters.
	\end{itemize}
	
	
	
	
	
	
	\section{Dataset}
	\label{sec:dataset}
	
	FRACTAL (Gaydon et al., 2024) was built by IGN from LiDAR HD data acquired in five regions of
	southern France. It contains 100,000 patches of $50 \times 50$~m, for a total of about 9.3 billion
	points (37 points per m$^2$ on average). Each point has its coordinates, the intensity of the returned
	signal, its return number and the number of returns of its laser pulse, and a colour (red, green, blue
	and near-infrared) taken from aerial images. The patches are split geographically into a training set
	(80,000 patches), a validation set (10,000) and a test set (10,000).
	
	For this project, the original LAZ files were converted to Parquet by the teacher and stored
	on Amazon S3, with one file per patch and one folder per split. The training set alone weighs about
	263~GB and the test set about 35~GB. Each row of a file is a point, with 21 columns; besides the
	attributes listed above, the other columns describe the acquisition (scan angle, GPS time, flight line)
	and are not used in this work.
	
	Every point is assigned to a class code.
	In the paper, these codes are grouped into seven classes (Table~\ref{tab:classes}). In our pipeline, the rare codes 65 and 66 (artefacts and virtual points) are assigned to \textit{other} and code 67 (miscellaneous buildings) to \textit{building}.
	
	\begin{table}[H]
		\centering
		\begin{tabular}{llrrr}
			\hline
			Class & LiDAR HD codes & Train (\%) & Val (\%) & Test (\%) \\
			\hline
			Other               & 1       & 0.56  & 0.53  & 0.66  \\
			Ground              & 2       & 39.0  & 39.1  & 40.5  \\
			Vegetation          & 3, 4, 5 & 57.0  & 56.9  & 54.1  \\
			Building            & 6       & 2.80  & 2.80  & 3.33  \\
			Water               & 9       & 0.52  & 0.49  & 1.20  \\
			Bridge              & 17      & 0.13  & 0.10  & 0.16  \\
			Permanent structure & 64      & 0.04  & 0.04  & 0.03  \\
			\hline
		\end{tabular}
		\caption{Classes and share of points per split (Gaydon et al., 2024).}
		\label{tab:classes}
	\end{table}
	
	The dataset is strongly imbalanced. Ground and vegetation represents about 96\% of the points, while
	bridges and permanent structures represent less than 0.2\%. 
	This matters for the evaluation, since a model that ignores the rare classes can still reach a high overall accuracy.
	It's why mIoU, that gives the same weight to every class is the main metric of the project



	\section{Tools}
	
	\section{Methodology}
	
	\subsection{Sampling}
	
	Each Parquet file corresponds to one patch, so a fraction of the dataset is defined as a fraction of
	the files of each split. Sampling files rather than points avoids reading the whole dataset before
	sampling it. For each split, the list of files is sorted, shuffled once with a fixed seed, and the
	first $k$ files are kept, where $k$ is the requested percentage of the files. As the list is shuffled
	only once, the samples are nested: the 1\% sample is included in the 2\% sample, which is included in
	the 5\% sample, and so on. Runs with different fractions are therefore comparable. Only the columns
	needed for the features and the label are read, which Parquet supports natively. The same fraction is
	applied to the training and test sets.
	
	\subsection{Preprocessing and Features}
	
	\subsection{Model}
	
	\subsection{Evaluation Metrics}
	
	\subsection{Cluster}
	
	\subsection{Executor Configuration}
	
	\subsection{Experiments}
	
	\section{Results}

	\subsection{Classification Quality}

	\subsection{Scalability}

	
	
	
\end{document}